% Part: first-order-logic
% Chapter: axiomatic-proofs
% Section: identity

\documentclass[../../../include/open-logic-section]{subfiles}

\begin{document}

\olfileid{fol}{axd}{ide}

\olsection{\usetoken{S}{identity}-সহ \usetoken{P}{derivation}}

!!{derivation}s-এ $\eq$ অন্তর্ভুক্ত করার জন্য শুধু নতুন কয়েকটি
স্বতঃসিদ্ধ-ছক যোগ করি। !!{derivation} ও~$\Proves$-এর সংজ্ঞা একই
থাকে; কেবল নতুন স্বতঃসিদ্ধগুলিকেও অনুমোদন করি।

\begin{defn}[!!{identity}-র স্বতঃসিদ্ধ]
\begin{align}
\ollabel{ax:id1} & \eq[t][t], \\
\ollabel{ax:id2} & \eq[t_1][t_2] \lif (!B(t_1) \lif
  !B(t_2)),
\end{align}
যেকোনো বদ্ধ পদ $t$, $t_1$, $t_2$-এর জন্য।
\end{defn}

\begin{prop}\ollabel{prop:sound}
  \olref{ax:id1} ও~\olref{ax:id2} স্বতঃসিদ্ধ দুটি বৈধ।
\end{prop}

\begin{proof}
  অনুশীলনী।
\end{proof}

\begin{prob}
\olref[fol][axd][ide]{prop:sound} প্রমাণ করো।
\end{prob}

\begin{prop}\ollabel{prop:iden1}
 যেকোনো বদ্ধ পদ~$t$ ও সেট~$\Gamma$-র জন্য $\Gamma \Proves \eq[t][t]$।
\end{prop}

\begin{prop}\ollabel{prop:iden2}
  বদ্ধ পদদুটির জন্য $\Gamma \Proves !A(t_1)$ এবং $\Gamma
  \Proves \eq[t_1][t_2]$ হলে $\Gamma \Proves !A(t_2)$।
  % BN-SRC-066 TARGET-CORRECTION-BEGIN
  \par\smallskip\noindent\textnormal{\small\textbf{উৎস-সংশোধন (BN-SRC-066):} হিমায়িত ইংরেজি স্বতঃসিদ্ধ-সংজ্ঞা $t,t_1,t_2$-কে বদ্ধ পদে সীমিত করে, কিন্তু পরের দুটি প্রতিজ্ঞার প্রথমটি “যেকোনো পদ” বলে এবং দ্বিতীয়টি কোনো সীমা জানায় না। প্রদত্ত স্বতঃসিদ্ধ-ছক থেকে ফলগুলি যাতে সত্যিই অনুসৃত হয়, বাংলা প্রতিজ্ঞা দুটিতে একই বদ্ধ-পদ পরিসর স্পষ্ট করা হয়েছে।} % BN-SRC-066 TARGET-CORRECTION-END
\end{prop}

\begin{proof}
নিচের !!{formula}-টি
\[
(\eq[t_1][t_2] \lif (!A(t_1) \lif !A(t_2)))
\]
হলো~\olref{ax:id2}-এর একটি রূপ। \MP দুবার প্রয়োগ করলে সিদ্ধান্তটি
পাওয়া যায়।
\end{proof}

\end{document}
